\section{Evaluation and observability}
\label{sec:eval}
\textbf{Harness.} \texttt{eval/run\_eval.py} sends each item to \texttt{/agent} (or \texttt{/chat} with a strategy),
one request at a time. It grades facts, behaviour, citation, retrieval of a gold clause and expected tools, and
records tokens, LLM calls, tool calls, latency and the final status. Multi-turn items are sent as consecutive turns
in one conversation. Failures get the first matching label: \emph{tool-call error} (a tool returned an error),
\emph{retrieval miss} (gold clauses exist but none were retrieved), \emph{hallucination} (wrong although the gold
clause was retrieved or the right tool ran), and \emph{reasoning error} (everything else).

\begin{table}[h]
\centering\small
\caption{End-to-end results. Cost is in tokens: the course endpoint is quota-based, not priced per token.}
\label{tab:results}
\setlength{\tabcolsep}{3.5pt}
\begin{tabular}{@{}lccccccccc@{}}
\toprule
Run & $n$ & Acc. & Behav. & Compl. & Cited & Recall & Tools & LLM & Tokens in / out \\
\midrule
Dev, agent & 50 & \multicolumn{8}{l}{\todoin{\texttt{run\_eval.py --mode agent} on gold\_v1}} \\
\textbf{Held-out, agent} & 40 & \multicolumn{8}{l}{\todoin{\texttt{--mode agent} on heldout\_v1 (Vijay's set), run once}} \\
Held-out, grounded prompt & 40 & \multicolumn{8}{l}{\todoin{single-prompt baseline run on heldout\_v1}} \\
\bottomrule
\end{tabular}
\end{table}

\todo{p50/p95 latency per setup; held-out results by category; failure taxonomy table (count and
one example per failure type: retrieval miss, hallucination, reasoning error, tool-call error); dev vs.\ held-out
gap and what it says about tuning on the dev set.}
\todo{critic on/off ablation (\texttt{max\_critic\_rounds=0}) on heldout\_v1: accuracy, critic
false-positive rate and token cost.}
\todo{LLM-as-judge faithfulness score (is every claim supported by the retrieved text?) next to the string-match
grade; not implemented.}

\textbf{Tracing.} Every answer's trace is built from LangGraph's own event stream and stored in Postgres
(\texttt{chat.trace\_steps}). Each row is one LLM, tool or retriever step with node, inputs (prompt messages, tool
arguments, search query), outputs (reply and tool calls, tool result, retrieved chunks with citations), tokens,
timings, error and a parent pointer, for example the retriever inside \texttt{search\_policy}. Each answer row links
to its trace through a \texttt{turn\_id}. The UI shows the trace next to the chat and on a trace page. Keeping the
trace in the application database keeps the deployment at three containers. It is what makes the failure taxonomy
possible: the trace shows the query the agent chose and what each retrieval stage returned, which separates
retrieval misses from generation errors.
